\section{Synthesis: the hypothesis, sentence by sentence}
\label{s5:sec:synthesis}

We now return to the hypothesis of \cref{s0:sec:question} and test each of its clauses against what
\cref{d1:sec:main,d2:sec:breaks,d3:sec:main,d4:sec:bfield} establish. The verdicts are summarised in \cref{s5:tab:verdict}. Throughout,
the relation $E=mc^2$ itself is inherited; what is at stake is only how the arena realises its three ingredients.

\subsection{``The energy is one undivided entity, and it is six-dimensional''}

This clause survives, in the field-strength reading of $\R^{3,3}$. Among quadratic forms on the six frame components, exactly one is
invariant under the diagonal rotations and under duality, and it is positive: $E_6=\tfrac12|z|^2=\tfrac12(|\bB|^2+|\bE|^2)$
(\cref{d1:sec:unique}; this restates Theorem~17.11(iii) of \cite{ES17}). Its symmetry group is $U(3)$, so the energy is ``unitary'' in a
literal sense, but only as a pointwise quadratic form: $U(3)$ is not a symmetry of Maxwell's equations on the frame, nor of the
split-form interior. In the closure of the force law the undivided energy of the \emph{momentum} field, $E_6(\pi)=\tfrac12|p|^2$, is a
Casimir, conserved by every Hamiltonian, together with its rest energy $E_0(\pi)$ and its Poynting vector $\mathbf S(\pi)$
(\cref{d3:sec:conserve}). The energy of the \emph{frame} field is not conserved. The other reading of ``six-dimensional'', three temporal
energies forming a vector in $V^+$, fails: splitting that vector lowers the interior mass, $m_4^2c^4=\mu^2c^4-|\mathcal E_\perp|^2$
(\cref{d1:sec:momentum}).

\subsection{``Mass emerges through an internal differentiation of that unity''}

This clause survives as an identity, not as a process. With no interaction and no second field, the undivided energy of a field point
splits as
\[
E_6^2=E_0^2+|\mathbf S|^2,\qquad E_0=\tfrac12|q|=\tfrac12|z\cdot z|,
\]
the complex Lagrange identity (\cref{d1:sec:identity}). In SI units it reads $u^2=c^2|\mathbf g|^2+u_0^2$, the field-point form of
$E^2=(pc)^2+(mc^2)^2$, with rest-mass density $u_0/c^2$. The rest part $E_0$ is the complete invariant of Lorentz transformations together
with duality, the Lorentz scalar of the stress tensor whose $00$ component is $E_6$ (Rainich), and the minimum of $E_6$ over boosts,
attained exactly in the rest frame (\cref{d1:sec:min}). All of this is inherited classical electrodynamics written in the arena's
variables. The arena adds which forms play which role: $E_6$ belongs to the duality class of Theorem~17.11 of \cite{ES17}, $E_0$ to the
Lorentz class, and $E=mc^2$ is the locus, the rest frame, where the two coincide.

The differentiation does not happen in time. In the closure the split $E_6^2=E_0^2+|\mathbf S|^2$ of the momentum field is fixed by the
initial data, because all three terms are Casimirs. What moves is the duality angle and the rotor rate, and the conserved total
$H=\tfrac12I_0\kappa^2+U(\theta)$ exchanges energy between a rate part and a Lorentz-invariant frame part (\cref{d3:sec:budget}). The mass
of a whole is not produced by interference either: two disjoint counter-propagating null pulses have $q\equiv0$ everywhere and $M=U>0$
(\cref{d1:sec:whole}).

\subsection{\texorpdfstring{``$c^2$ is the symmetric part''}{"c\textasciicircum 2 is the symmetric part"}}

This clause survives in one precise form. The constant $c$ is the unique isotropic conversion between the temporal and the spatial triple,
forced by Schur's lemma (Theorem~3.2 of \cite{ES17}). It is a conversion, not a symmetry, and its weight in $E_6$ is a unit convention
(\cref{d1:sec:c}). The form in which the clause is exact is the decomposition of a shift operator into its two summands
(Theorem~17.3 of \cite{ES17}):
\[
\End(V)=\so{3,3}\oplus\{X:X^\dagger=X\},\qquad \text{shift}=b+\kappa G,\qquad (\kappa G)^2=-\frac{c^2}{\ell^2}\,\id .
\]
The $\eta$-self-adjoint summand is the symmetric part: it is the summand in which factorisation holds, the one that binds in neutral
signature, and the one that carries $c$, through $\kappa=c/\ell$. The antisymmetric summand is a connection. Reading the first as ``$c^2$''
and the second as ``$b$'' is interpretive, but the decomposition, and the fact that only the first summand binds, are fixed by the paper.

\subsection{\texorpdfstring{``$m$ is what breaks the symmetry''}{"m is what breaks the symmetry"}}

This clause is true for the right symmetries and false for the ones it is usually read about (\cref{d2:sec:breaks}).
\begin{enumerate}
\item A nonzero rest energy breaks duality. The duality $U(1)$ generated by $G$ acts on the Lorentz invariant by $q\mapsto e^{2i\theta}q$, so
a field with $E_0\neq0$ keeps only $\{\pm\id\}$, and $\arg q$ labels the broken directions. At a null field point, duality joins the
little group, which becomes $\mathfrak e(2)$ with helicity as its compact generator. This is inherited, pulled back to the frame.
\item For Dirac matter the mass breaks the chiral and Pauli--G\"ursey symmetries, $\uu1_\Omega\oplus\su2\to\uu1_J$, with three Goldstone
directions, one of them the pseudoscalar mass angle (Theorem~18.5 of \cite{ES17}, re-verified).
\item The mass breaks no spacetime symmetry, not even the six-dimensional one. The Dirac mass form $T=e_1e_2e_3$ is invariant under all of
$\mathrm{Spin}(3,3)$ and is the only such kernel (Corollary~18.4 of \cite{ES17}). Within $\so{3,3}$ it is the kinetic term, where $c$
enters, that cuts the symmetry down to the interior Lorentz algebra and $J$. In this sense the hypothesis is inverted.
\item In the closure the rest energy of the momentum field is exactly the amplitude of the potential for the duality angle,
$R_0=2|\nu_1+i\nu_2|\,E_0(\pi)$, with $(\nu_1,\nu_2)$ the Lorentz-class inertia coefficients of
\cref{d1:sec:main}. A null (massless) momentum field leaves the angle flat and the rate constant; a massive one makes it a pendulum
whose squared small-oscillation frequency is
$4|\nu_1+i\nu_2|E_0(\pi)/(mI_0)$ (\cref{d3:sec:pendulum}). This restates Theorem~17.13 and
Corollary~17.14 of \cite{ES17}. Reading the angle as a pseudo-Goldstone boson of duality is interpretive: the breaking is explicit, carried
by the Lorentz-class inertia, and every pendulum has this shape.
\end{enumerate}

\subsection{\texorpdfstring{``$m$ is generated from the $b$-particle of Yang--Mills''}{"m is generated from the b-particle of Yang-Mills"}}

Classically, this clause fails, and the failure is sharp (\cref{d4:sec:bfield}).
\begin{enumerate}
\item The binding rate is never a gauge potential. For every connection $b$ the field equation factorises covariantly,
$(D-\kappa G)(D+\kappa G)X=0$ with $D=\partial_\tau+b$; $\kappa G$ transforms as a tensor, $b$ as a connection. Without $\kappa$, a constant $b$
is pure gauge along the worldline: no rest frequency, no mass scale.
\item An admissible Yang--Mills field is compact, since no gauge boson is a boost (Corollary~2.4 of \cite{ES17}). Its part commuting with
$G$ splits the rest frequency into charge-labelled lines $\kappa$, $\kappa\pm\beta$, removable along a single worldline by a rotating frame
(Larmor's theorem). This is a spectrum without a scale, since $\beta/\kappa$ is free. A purely anticommuting part produces Jordan blocks and
secular growth. A mixed connection keeps every orbit bounded exactly when it rotates the temporal and the spatial triple at equal rates about
axes that are not opposite, outside one exceptional ratio of the rates to $\kappa$ (\cref{d4:thm:screen}).
\item At three constituents the part of the Pati--Salam algebra commuting with the binding root is $\so4_\Delta\cong\su2\oplus\su2$
(Theorem~17.21 of \cite{ES17}): two copies of Yang and Mills' isospin algebra. It splits the rest frequency by $\su2$ weights and supplies
no scale.
\item In the selected vacuum the order is the reverse of the hypothesis. The Pati--Salam vector bosons acquire mass by eating nine
Goldstone directions (Theorem~22.2 of \cite{ES17}), rather than generating the mass of anything.
\item At the quantum level, one-gauge-boson exchange favours the Dirac singlet channel $(\mathbf1,\mathbf2,\mathbf2)$ and the
colour-breaking $(\mathbf6,\mathbf1,\mathbf1)$, and disfavours both the $(\mathbf{15},\mathbf2,\mathbf2)$ needed for the Georgi--Jarlskog factor
and the $(\mathbf{10},\mathbf1,\mathbf3)$ that the paper's vacuum uses. This is inherited group theory applied to the arena's representations.
\end{enumerate}
The one role left open for the $b$-field is the scale. The algebra fixes $c$ and leaves $\ell$, hence $\kappa$, free; classical Yang--Mills
theory in four dimensions has no scale either, but its quantum theory generates one by dimensional transmutation,
$\Lambda_b=\mu\exp[-8\pi^2/(\beta_0g^2(\mu))]$. The identification $\hbar\kappa=\Lambda_b$, that is $\ell=\hbar c/\Lambda_b$, is a conjecture,
not derived. Even granted, reading the rate as a mass, $mc^2=\hbar\kappa$, is a further assumption: in the interior dynamics $m$ and $\kappa$
are independent inputs, the spectrum does not depend on $m$, and the classical rescaling invariant is $m\kappa$.

\subsection{Particles, antiparticles, the vacuum and the Goldstone bosons}

The original questions of this investigation have the following answers in the arena (\cref{d2:sec:senses,d2:sec:vacua,d2:sec:goldstone}).
For Dirac matter at rest the flow is a rotation generated by the six-dimensional mass form itself, $\Psi(\lambda)=e^{-m\lambda T}\Psi_0$, with
rate equal to the mass. Particle and antiparticle are its two senses, and charge conjugation, the surplus plane, exchanges them. The binding
rotation of the frame also has two senses, but parity exchanges them, not charge conjugation, so its rate $\kappa=c/\ell$ is a conversion rate
and its two senses are not a particle--antiparticle pair. The Fock vacuum of one momentum is unique and neutral, with energy $-2E$
($-2mc^2$ at rest). The frame has no ground state, its spectrum being $\hbar\kappa\,\Z$, but its zero-point energy cancels in neutral
signature, and this cancellation survives, frequency by frequency, every connection whose motion stays bounded (\cref{d4:sec:bfield}).
Goldstone directions: three at one constituent in the Dirac phase; nine at three constituents, all eaten, with no Majoron because
$[J,\Phi_0]=[J\Gamma,\Phi_0]=2[t^R_3,\Phi_0]$; and one axion-like direction of charge $1-\Gamma$ at the electroweak stage, which a realistic
potential must lift explicitly.

\begin{table}[t]
\centering
\small
\begin{tabular}{@{}>{\raggedright\arraybackslash}p{3.6cm}>{\raggedright\arraybackslash}p{6.6cm}>{\raggedright\arraybackslash}p{4.0cm}@{}}
\toprule
clause & what the arena gives & verdict\\
\midrule
energy undivided, six-dimensional & $E_6=\tfrac12(|\bB|^2+|\bE|^2)$, unique, positive, symmetry $U(3)$ pointwise; Casimir of the momentum field &
holds (inherited / restates \cite{ES17})\\
mass from internal differentiation & $E_6^2=E_0^2+|\mathbf S|^2$, $E_0=\min_{\rm boosts}E_6$; split fixed by initial data in the closure &
holds as an identity, not as a process\\
$c^2$ the symmetric part & $c$ a Schur conversion; $\kappa G$ the $\eta$-self-adjoint summand of the shift, the only one that binds &
holds in the shift decomposition (interpretive reading of a fixed result)\\
$m$ breaks the symmetry & breaks duality and the chiral algebra; breaks neither Lorentz nor $\mathrm{Spin}(3,3)$ & true for internal symmetries,
false for spacetime ones\\
$m$ generated from the $b$-field & $b$ is the connection summand: splits or tears the binding, never creates it; vector bosons are massive by
eating Goldstones & fails classically; scale by transmutation open\\
\bottomrule
\end{tabular}
\caption{The hypothesis against the arena.}
\label{s5:tab:verdict}
\end{table}

\subsection{The hypothesis, restated so that the algebra supports it}

Collecting the verdicts, the strongest statement of the hypothesis that the arena supports is the following. \emph{The undivided energy
is the duality-invariant, unitary energy $\tfrac12(|\bB|^2+|\bE|^2)$ of the six field components. The constant $c$ is the unique isotropic
conversion between times and lengths, and it enters the dynamics through the $\eta$-self-adjoint part of the shift, the binding root. The
mass is the Lorentz-and-duality invariant $\tfrac12|q|$, the minimum of the energy over boosts, and $E=mc^2$ is the rest frame, where the
two coincide. What the mass breaks is duality and, for matter, chirality, never Lorentz invariance or $\mathrm{Spin}(3,3)$. The Yang--Mills
$b$-field is the antisymmetric part of the same shift: it can split or destabilise the binding, but it cannot create it, and the only
role left for it in generating mass is to supply, by dimensional transmutation, the one scale $\ell$ that the algebra leaves free.}

\subsection{Open problems}

\begin{enumerate}
\item A mechanism fixing $\ell$, and with it $\kappa$, from the arena or from the quantum dynamics of a $b$-field, so that $\hbar\kappa=\Lambda_b$
becomes a derivation; and a reason, absent so far, to identify the interior rate with a particle mass.
\item A dynamical process in which rest energy is produced rather than carried: in the closure the rest energy of the momentum field is a
Casimir, while that of the frame field drifts without bound.
\item A common Lorentz action for the two readings of $\R^{3,3}$: the field-strength reading, where $E_6$ and $E_0$ live, and the interior
reading, where the Dirac mass form lives, exclude each other on the same six directions (Proposition~12.6 of \cite{ES17}).
\item The classification of compact connections that keep the binding bounded, beyond the commuting ones and the equal-rate family found here.
\end{enumerate}
